
\medskip

%On dispose des informations suivantes :
%
%\medskip
%
%Toutes les valeurs présentes sur les schémas sont en millimètres.
%
%\begin{center}
%\begin{tabularx}{\linewidth}{>{\centering \arraybackslash}X m{2cm}m{5cm}}
%Dimensions de la remorque 	&	&Longueur du fusil sous-marin\\
%\psset{unit=0.8cm}
%\begin{pspicture}(0,-0.5)(7.5,4)
%\psframe[fillstyle=solid,fillcolor=lightgray](3,0.5)(7.2,3)
%\psline[linewidth=3pt](3,1.75)(0,1.75)
%\psline{<->}(7.4,0.5)(7.4,3)
%\psline{<->}(3,0.2)(7.2,0.2)
%\rput{90}(7.6,1.75){\np{1350}}
%\uput[d](5.1,0.2){\np{1800}}
%\psline{<->}(0,-0.4)(7.2,-0.4) 
%\uput[d](3.6,-0.4){\np{3000}}
%\end{pspicture}				&	&\psscalebox{0.35}{\includegraphics{fusil}}
%
%\np{2100}\\
%\end{tabularx}
%\end{center}
%
%On suppose que le fond de la remorque est un rectangle.
%
%Le fusil sous-marin peut-il être placé \og à plat \fg{} dans la remorque ? Justifier la réponse.
Dans le rectangle la plus grande longueur est celle d'une diagonale ou encore l'hypoténuse d'un triangle rectangle de côtés \np{1800} et \np{1350}.

D'après le théorème de Pythagore cette diagonale $d$ vérifie :

$d^2  = \np{1800}^2 + \np{1350}^2 = \np{3240000} + \np{1822500} = \np{5062500}$.

Donc $d = \sqrt{\np{5062500}} = \np{2250} > \np{2100}$.

Donc s'il n'est pas trop large le fusil pourra être placé à plat au fond de la remorque.
\medskip

%\emph{Dans cet exercice, toute trace de recherche, même incomplète ou non fructueuse, sera prise en compte dans l'évaluation.}



